\section{Overview}\label{sec:intro}

Tigers \& Goats is a two-player board game on 23 squares. Three tigers start on the board. The goats enter
one at a time until 15 have been placed, and then slide. A tiger captures a goat by jumping over it; six
captures win for the tigers. The goats win by leaving the tigers without a legal move.

The project studies the game in two ways, and this document gives the mathematics of both.

\begin{itemize}
\item \textbf{Learning (Part II).} A neural network plays one side against a fixed opponent and improves by
  trial and error. Section~\ref{sec:fp} derives the method from first principles. Sections
  \ref{old:3}--\ref{old:6} then state exactly what each saved model optimized.
\item \textbf{Solving (Part III).} Dynamic programming computes the true value of every position, with no
  learning and no approximation. Section~\ref{sec:dp} gives the recursion, the algorithms and the proof that
  the tables are right. Section~\ref{sec:longest} finds the longest forced wins and proves them.
\end{itemize}

The two meet in Part IV. Because every position's exact value is known, every move a model makes can be
judged as value-preserving or as a blunder. Section~\ref{sec:comparisons} does this for all models.

\paragraph{Main results.}
\begin{itemize}
\item The game is a draw with perfect play, whichever side moves first.
\item The longest forced win a real game can reach is 213 plies when the goats move first and 214 when the
  tigers do. Both are goat wins. The tigers' longest forced win is 67 plies.
\item No trained model beats a perfect player. Against a perfect tiger that looks for traps, the goat
  models lose between 12\% and 98\% of their games. Against the matching perfect goat, all three tiger
  models lose every game.
\item A perfect player is not the best exploiter of a weak opponent. The models win more often against the
  scripted tigers they were trained on than the perfect goat does.
\item Masking teaches a network nothing about which moves are legal. A supervised legality loss does.
\end{itemize}

\paragraph{Sources.} Parts I, II and IV extend the earlier \code{docs/TRAINING\_MATHEMATICS.md} (v1.2), which this
document replaces. Line numbers of the form \code{env\_tng.py:399} refer to commit \code{11ee06e} (v1.1) of that
file. Two later changes to it are noted where they matter: the option \code{terminate\_on\_invalid} (v1.3) and a
correction to one capture (Section~\ref{sec:rulefix}).
